\section{Parity-graded sectors}%
\label{sec:asymex_parity}

Consider a two-state physical site with basis $\{\ket{0},\ket{1}\}$,
where $1$ denotes occupation. Define the bond-two MPS tensor by
\begin{align}
	A^0=P=\begin{pmatrix}1&0\\0&0\end{pmatrix},
	\qquad
	A^1=X=\begin{pmatrix}0&1\\1&0\end{pmatrix}.
	\label{eq:asymex_parity_letters}
\end{align}
For $N\geq1$, write its periodic amplitudes as
\begin{align}
	a_N(x)=\tr(A^{x_0}\cdots A^{x_{N-1}})
	\qquad (x\in\{0,1\}^N).
	\label{eq:asymex_parity_amplitudes}
\end{align}
The associated unnormalized state is
\begin{align}
	\ket{V^{(N)}(A)}
	=\sum_{x\in\{0,1\}^N}a_N(x)\ket{x},
	\label{eq:asymex_parity_state}
\end{align}
where $\ket{x}=\ket{x_0}\otimes\cdots\otimes\ket{x_{N-1}}$.
As in the elementary MPS examples of
Ref.~\cite[Appendix~A]{Cirac2021Matrix}, the state can be understood
directly from its tensor letters. The matrix $P$ selects virtual
coordinate $0$, while $X$ exchanges the two virtual coordinates.

\subsection{Occupation parity and periodic amplitudes}

Let $z=\diag(1,-1)$ act on the physical site, and let
$Z=\diag(1,-1)$ act on the bond space. Although their matrices are
equal, these operators act on different spaces. If
$n_1(x)=\sum_{v=0}^{N-1}x_v$, the physical parity operator has basis action
\begin{align}
	z^{\otimes N}\ket{x}=(-1)^{n_1(x)}\ket{x}.
	\label{eq:asymex_parity_basis}
\end{align}

\begin{lemma}[Occupation parity]
	\label{lem:asymex_parity_occupation}
	For each physical label $i\in\{0,1\}$, the tensor satisfies
	\begin{align}
		ZA^iZ=(-1)^iA^i.
		\label{eq:asymex_parity_covariance}
	\end{align}
	Consequently, $a_N(x)=0$ whenever $n_1(x)$ is odd, and
	\begin{align}
		z^{\otimes N}\ket{V^{(N)}(A)}=\ket{V^{(N)}(A)}.
		\label{eq:asymex_parity_invariance}
	\end{align}
\end{lemma}

\begin{proof}
	Direct multiplication gives $ZPZ=P$ and $ZXZ=-X$.
	Since $Z^2=\Id_2$, multiplying~(\ref{eq:asymex_parity_covariance})
	along the word yields
	\begin{align}
		Z(A^{x_0}\cdots A^{x_{N-1}})Z
			                            =(-1)^{n_1(x)}A^{x_0}\cdots A^{x_{N-1}}.
		\notag
	\end{align}
	Taking the trace and using cyclicity gives
	$a_N(x)=(-1)^{n_1(x)}a_N(x)$. Thus odd occupation forces
	$a_N(x)=0$. Applying~(\ref{eq:asymex_parity_basis}) to the
	expansion~(\ref{eq:asymex_parity_state}) proves
	(\ref{eq:asymex_parity_invariance}).
\end{proof}

Occupation parity does not completely determine the support of the
state. The positions of the zeros impose an additional restriction.

\begin{lemma}[Amplitudes from cyclic runs]
	\label{lem:asymex_parity_runs}
	Suppose $x$ contains $m\geq1$ zeros. Read the sites cyclically, and
	let $r_1,\ldots,r_m$ be the numbers of ones between successive
	zeros. Define $b(r)=1$ for even $r$ and $b(r)=0$ for odd $r$.
	Then
	\begin{align}
		a_N(x)=\prod_{j=1}^{m}b(r_j).
		\label{eq:asymex_parity_runs}
	\end{align}
	For the word with no zeros,
	\begin{align}
		a_N(1,\ldots,1)=1+(-1)^N.
		\label{eq:asymex_parity_all_one}
	\end{align}
\end{lemma}

\begin{proof}
	The identities $X^2=\Id_2$ and $PXP=0$ imply
	$PX^rP=b(r)P$ and $\tr(PX^r)=b(r)$. Rotate the word to
	start at a zero, using cyclicity of the trace. Successive
	contractions give
	\begin{align}
		a_N(x)
		 & =\tr(PX^{r_1}PX^{r_2}\cdots PX^{r_m}) \notag       \\
		 & =\left(\prod_{j=1}^{m-1}b(r_j)\right)\tr(PX^{r_m})
		=\prod_{j=1}^{m}b(r_j).
		\notag
	\end{align}
	For the all-one word, $X^N=\Id_2$ when $N$ is even and
	$X^N=X$ when $N$ is odd. Their traces are $2$ and $0$,
	respectively, which proves~(\ref{eq:asymex_parity_all_one}).
\end{proof}

Thus a word containing a zero has amplitude one exactly when every
cyclic run of ones has even length. For example, $0101$ has even
occupation but zero amplitude, since both runs have length one.
The all-zero word has amplitude one at every positive length, so
$\ket{V^{(N)}(A)}$ is always nonzero. These statements describe
occupation parity in the specified computational basis. No
Hamiltonian or operators satisfying the canonical anticommutation
relations have been specified, so they do not identify this state
with a Kitaev-chain state.

\begin{theorem}[Normality of the graded block]
	\label{thm:asymex_parity_normal}
	\lean{ParityGraded.parA_isNormal}
	\leanok
	\uses{def:asymex_parity, def:normal}
	The block $A$ is normal: its length-two words span $\MN{2}$.
\end{theorem}

\begin{proof}
	\leanok
	Let $E_{ab}$ denote the matrix unit with its sole nonzero entry
	equal to one in row $a$, column $b$. Direct multiplication gives
	$E_{00}=A^0A^0$, $E_{01}=A^0A^1$, and $E_{10}=A^1A^0$.
	The remaining unit is $E_{11}=A^1A^1-A^0A^0$, so all four
	matrix units lie in the length-two word span.
\end{proof}

This is word normality, established by spanning the virtual matrix
algebra. It makes no assertion that the transfer map has been
normalized to spectral radius one.

\subsection{Sign-weighted blocks and the five-dimensional source}

Multiplying both letters of $A$ by $-1$ gives one minus sign per
site, not one minus sign per occupied site. For a word
$w=(i_1,\ldots,i_{|w|})$, write
$A^w=A^{i_1}\cdots A^{i_{|w|}}$, with the identity matrix as the
empty product. Then
\begin{align}
	(-A)^w=(-1)^{|w|}A^w.
	\label{eq:asymex_parity_weight}
\end{align}
The two weights $+1$ and $-1$ therefore produce a factor
$1+(-1)^N$ in the sum of their periodic amplitudes. The following
five-dimensional tensor realizes that sum without being a direct
sum in its displayed coordinates.

\begin{definition}[The graded pair and the five-dimensional source]
	\label{def:asymex_parity}
	\lean{ParityGraded.parA, ParityGraded.parTargets, ParityGraded.parB}
	\leanok
	\uses{def:mps_tensor}
	Let $A$ be the tensor in~(\ref{eq:asymex_parity_letters}).
	The two target tensors are $C_+=A$ and $C_-=-A$.
	The source is the bond-dimension-five tensor with letters
	\begin{align}
		B^0
		 & =\begin{pmatrix}
			    -1 & 1 & -1 & 1  & 3 \\
			    2  & 2 & -2 & -1 & 0 \\
			    1  & 2 & -2 & 0  & 2 \\
			    1  & 3 & -3 & -1 & 2 \\
			    -1 & 0 & 0  & 1  & 2
		    \end{pmatrix}
		\label{eq:asymex_parity_source_zero}
	\end{align}
	and
	\begin{align}
		B^1
		 & =\begin{pmatrix}
			    3  & 1 & 0  & -2 & -3 \\
			    -3 & 1 & -1 & 2  & 5  \\
			    -1 & 2 & -1 & 1  & 3  \\
			    0  & 1 & -1 & -1 & 1  \\
			    2  & 1 & 0  & -1 & -2
		    \end{pmatrix}.
		\label{eq:asymex_parity_source_one}
	\end{align}
\end{definition}

Neither source letter has a zero row or a zero column. Nevertheless,
a simultaneous change of bond coordinates reveals diagonal blocks
$A$, $-A$, and a one-dimensional zero block. The latter contributes
nothing to positive-length traces.

\begin{center}
	% Formula: the periodic trace tr(B^w)=(1+(-1)^{|w|})tr(A^w) of
	% Theorem~\ref{thm:asymex_parity_trace}; the source is not presented
	% as a literal stack of $A$ and $-A$, so the ring carries $B$ with
	% its coefficient, not a two-row stack.
	% Source: Definition~\ref{def:asymex_parity}.
	% Ink-to-index map: each node is the bond-dimension-five source
	% tensor $B$; the upper legs carry the letters of the word.
	% Contraction set: the horizontal virtual legs, closed around the
	% ring by the trace; no leg is left open.
	% Boundary signature: (virtual-west=0 | virtual-east=0 |
	% physical-up=3 | physical-down=0) for the finite schematic shown.
	% The physical legs remain open for the state; fixing their labels
	% gives the scalar coefficient tr(B^w). No virtual leg remains open.
	\tenkzkernel
	\begin{tenkz}[cols=4, boundary=periodic, physical=up]
		\tn[label pos=135, ports={90:physical:$i_1$}]{B} &
		\tn[label pos=135]{B} & \tn[skin=dots]{} &
		\tn[label pos=135, ports={90:physical:$i_N$}]{B}
	\end{tenkz}
	\qquad $\tr(B^w)$, computed in
	Theorem~\ref{thm:asymex_parity_trace}
\end{center}

\subsection{Asymmetric compression and length cancellation}

\begin{theorem}[Compression of the graded pair]
	\label{thm:asymex_parity_compression}
	% Principal declaration: ParityGraded.parityGraded_isReduction.
	\lean{ParityGraded.parityGraded_compression,
		ParityGraded.parityGraded_isReduction,
		ParityGraded.parityGraded_left_mul_right_of_ne,
		ParityGraded.parityGraded_dim_eq,
		ParityGraded.parityGraded_evalWord_remainder_eq_zero,
		ParityGraded.parityGraded_remainder_sq_ne_zero}
	\leanok
	\uses{def:asymex_parity, thm:asym_multiblock}
	The source of Definition~\ref{def:asymex_parity} carries the block
	triangular form of Theorem~\ref{thm:asym_multiblock} for
	$S=\{+,-\}$, with targets $C_+=A$, $C_-=-A$, one zero slot
	$z=1$, and dimension count $5=2+2+1$.
	There are compression maps $V_s\in\C^{5\times2}$ and
	$W_s\in\C^{2\times5}$ such that, for every $s,t\in S$,
	\begin{align}
		W_sV_t=\delta_{st}\Id_2.
		\label{eq:asymex_parity_biorthogonal}
	\end{align}
	For every word $w$ and each $s\in S$,
	\begin{align}
		W_sB^wV_s=C_s^w.
		\label{eq:asymex_parity_reduction}
	\end{align}
	Define the remainder letters by
	\begin{align}
		R^i=B^i-V_+A^iW_+-V_-(-A^i)W_-.
		\label{eq:asymex_parity_remainder}
	\end{align}
	Every product $R^w$ with $|w|\geq3$ vanishes, whereas
	$R^0R^0\neq0$. Thus the bound $|S|+z=3$ in clause (vi) of
	Theorem~\ref{thm:asym_multiblock} is attained.
\end{theorem}

\begin{proof}
	\leanok
	\uses{thm:asymex_explicit_gauge}
	Use the integer change of coordinates
	\begin{align}
		G=\begin{pmatrix}
			  0  & -1 & 1  & 0  & 0  \\
			  1  & 1  & 0  & -1 & -1 \\
			  -1 & 0  & 0  & 1  & 1  \\
			  1  & 0  & 0  & 0  & -1 \\
			  0  & 1  & -1 & 0  & 1
		  \end{pmatrix}.
		\label{eq:asymex_parity_gauge}
	\end{align}
	Its inverse is
	\begin{align}
		G^{-1}=\begin{pmatrix}
			       1 & 0 & 0 & 1 & 1 \\
			       0 & 1 & 1 & 0 & 0 \\
			       1 & 1 & 1 & 0 & 0 \\
			       0 & 0 & 1 & 1 & 0 \\
			       1 & 0 & 0 & 0 & 1
		       \end{pmatrix}.
		\label{eq:asymex_parity_gauge_inv}
	\end{align}
	Direct multiplication verifies both inverse identities and gives
	\begin{align}
		GB^0G^{-1}
		=\begin{pmatrix}
			 1 & 0 & 1  & 0 & 1 \\
			 0 & 0 & 0  & 1 & 0 \\
			 0 & 0 & -1 & 0 & 2 \\
			 0 & 0 & 0  & 0 & 1 \\
			 0 & 0 & 0  & 0 & 0
		 \end{pmatrix}.
		\label{eq:asymex_parity_conj_zero}
	\end{align}
	For the second letter,
	\begin{align}
		GB^1G^{-1}
		=\begin{pmatrix}
			 0 & 1 & 0  & 1  & 0 \\
			 1 & 0 & 2  & 0  & 1 \\
			 0 & 0 & 0  & -1 & 1 \\
			 0 & 0 & -1 & 0  & 0 \\
			 0 & 0 & 0  & 0  & 0
		 \end{pmatrix}.
		\label{eq:asymex_parity_conj_one}
	\end{align}
	With block sizes $(2,2,1)$, the matrices in
	(\ref{eq:asymex_parity_conj_zero}) and
	(\ref{eq:asymex_parity_conj_one}) are block upper triangular,
	with diagonal blocks $A^i$, $-A^i$, and $0$.
	As in Theorem~\ref{thm:asymex_explicit_gauge}, take $W_+$ and
	$W_-$ to be rows $(0,1)$ and $(2,3)$ of $G$, respectively,
	and $V_+$ and $V_-$ to be the corresponding columns of $G^{-1}$.
	The inverse identities give~(\ref{eq:asymex_parity_biorthogonal}).
	Diagonal blocks multiply under word evaluation, giving
	(\ref{eq:asymex_parity_reduction}).

	Subtracting the two diagonal target blocks leaves a strictly
	upper triangular three-block remainder in these coordinates.
	The compression remainder bound therefore gives $R^w=0$
	for $|w|\geq3$. To check sharpness in the original coordinates,
	substitution into~(\ref{eq:asymex_parity_remainder}) gives
	\begin{align}
		R^0=\begin{pmatrix}
			    -1 & 2 & -2 & 1 & 3 \\
			    1  & 2 & -2 & 0 & 1 \\
			    0  & 3 & -3 & 1 & 3 \\
			    0  & 3 & -3 & 0 & 3 \\
			    -1 & 1 & -1 & 1 & 2
		    \end{pmatrix}.
		\notag
	\end{align}
	Its square has $(0,1)$ entry $2$, so $R^0R^0\neq0$.
\end{proof}

\begin{theorem}[Periodic trace of the sign-weighted pair]
	\label{thm:asymex_parity_trace}
	\lean{ParityGraded.parityGraded_trace_evalWord}
	\leanok
	\uses{def:asymex_parity}
	For every nonempty word $w$,
	\begin{align}
		\tr(B^w)=(1+(-1)^{|w|})\tr(A^w).
		\label{eq:asymex_parity_trace}
	\end{align}
\end{theorem}

\begin{proof}
	\leanok
	\uses{thm:asymex_parity_compression, prop:asym_compression_trace}
	Proposition~\ref{prop:asym_compression_trace} applied to the datum of
	Theorem~\ref{thm:asymex_parity_compression} gives
	$\tr(B^w)=\tr(A^w)+\tr((-A)^w)$. By
	(\ref{eq:asymex_parity_weight}), the second summand is
	$(-1)^{|w|}\tr(A^w)$, proving~(\ref{eq:asymex_parity_trace}).
\end{proof}

Define $\ket{V^{(N)}(B)}$ by the same periodic contraction as
in~(\ref{eq:asymex_parity_state}), with $A$ replaced by $B$.
Applying~(\ref{eq:asymex_parity_trace}) to each basis coefficient
gives, for $N\geq1$,
\begin{align}
	\ket{V^{(N)}(B)}
	=(1+(-1)^N)\ket{V^{(N)}(A)}.
	\label{eq:asymex_parity_length}
\end{align}
Thus the source produces the zero vector at odd length and twice
the unnormalized $A$ state at even length. This is a length filter,
not an occupation-parity projection on a fixed physical space.
At odd length, for instance, the all-zero amplitude of $A$ is one,
but its amplitude in $B$ vanishes. At even length, every amplitude
is doubled, and the run restrictions of
Lemma~\ref{lem:asymex_parity_runs} remain unchanged.

The positive-length hypothesis is necessary. The empty-word
traces are $\tr(\Id_5)=5$ and $\tr(\Id_2)=2$, whereas substituting
$|w|=0$ into~(\ref{eq:asymex_parity_trace}) would give $5=4$.
The missing contribution is the empty trace of the
one-dimensional zero block, which is one; its traces vanish only
at positive length.

\begin{theorem}[The weight $-1$ block has no sitewise intertwiner]
	\label{thm:asymex_parity_no_intertwiner}
	\lean{ParityGraded.parNeg_right_intertwiner_eq_zero,
		ParityGraded.parNeg_left_intertwiner_eq_zero}
	\leanok
	\uses{def:asymex_parity}
	If $V\in\C^{5\times2}$ satisfies $B^iV=V(-A^i)$ for both letters, then
	$V=0$; if $W\in\C^{2\times5}$ satisfies $WB^i=(-A^i)W$ for both
	letters, then $W=0$.
\end{theorem}

\begin{proof}
	\leanok
	Each condition is a homogeneous linear system in ten unknown
	entries. For the right intertwiner, select entries
	$(3,0)$, $(4,0)$, $(2,1)$, $(3,1)$, and $(4,1)$ of
	$B^0V=V(-A^0)$, and all five entries of column $0$ of
	$B^1V=V(-A^1)$. Rational linear combinations of these ten
	equations force all entries of column $0$ of $V$ to vanish.
	The selected equations for letter $1$ then force column $1$
	to vanish.

	For the left intertwiner, select entries $(0,4)$, $(1,0)$,
	$(1,2)$, $(1,3)$, and $(1,4)$ of $WB^0=(-A^0)W$, and all
	five entries of row $0$ of $WB^1=(-A^1)W$. Rational linear
	combinations force row $0$ of $W$ to vanish, after which the
	letter-$1$ equations force row $1$ to vanish.
\end{proof}

The nonzero maps $W_-$ and $V_-$ in
(\ref{eq:asymex_parity_reduction}) therefore give a word-level
compression onto $-A$ without either map being a sitewise
intertwiner. The trace identity and the sharp remainder bound
hold despite this obstruction.
